\documentclass[10pt,a4paper]{article}

\usepackage[utf8]{inputenc}
\usepackage[T1]{fontenc}
\usepackage[french]{babel}
\usepackage{amsmath,amssymb}
\usepackage{geometry}
\usepackage{enumitem}
\usepackage{multicol}
\usepackage{xcolor}
\usepackage{mdframed}
\usepackage{lmodern}
\usepackage{microtype}
\usepackage{fancyhdr}
\usepackage{titlesec}
\usepackage{parskip}

\geometry{top=1.8cm, bottom=2cm, left=1.8cm, right=1.8cm}

\pagestyle{fancy}
\fancyhf{}
\fancyhead[L]{\small\textit{Mathématiques, fiche de révision}}
\fancyhead[R]{\small\textit{Opérateurs vectoriels}}
\fancyfoot[C]{\small\thepage}
\fancyfoot[R]{\tiny\textcopyright\ Rohan Fossé}
\renewcommand{\headrulewidth}{0.3pt}

\definecolor{bleu}{RGB}{30,60,120}
\definecolor{vert}{RGB}{0,110,60}
\definecolor{rouge}{RGB}{160,20,20}
\definecolor{gris}{gray}{0.93}
\definecolor{jauneclair}{RGB}{255,250,220}
\definecolor{vertclair}{RGB}{220,245,220}
\definecolor{bleutresleger}{RGB}{230,240,255}

\newmdenv[backgroundcolor=gris, linewidth=0pt,
  innerleftmargin=8pt, innerrightmargin=8pt,
  innertopmargin=5pt, innerbottommargin=5pt,
  skipabove=4pt, skipbelow=4pt]{grisbox}

\newmdenv[backgroundcolor=jauneclair, linecolor=rouge, linewidth=1pt,
  innerleftmargin=8pt, innerrightmargin=8pt,
  innertopmargin=5pt, innerbottommargin=5pt,
  skipabove=4pt, skipbelow=4pt]{alertbox}

\newmdenv[backgroundcolor=bleutresleger, linecolor=bleu, linewidth=1pt,
  innerleftmargin=8pt, innerrightmargin=8pt,
  innertopmargin=5pt, innerbottommargin=5pt,
  skipabove=4pt, skipbelow=4pt]{theobox}

\titleformat{\section}[block]{\normalsize\bfseries\color{bleu}}{}{0pt}{}[\vspace{-3pt}\rule{\linewidth}{0.5pt}]
\titlespacing{\section}{0pt}{10pt}{4pt}

\titleformat{\subsection}[runin]{\small\bfseries\color{vert}}{}{0pt}{}[\ :]
\titlespacing{\subsection}{0pt}{4pt}{0pt}

\begin{document}

\begin{center}
  {\Large\bfseries\color{bleu} Fiche de révision : Opérateurs vectoriels}\\[3pt]
  {\small Gradient · Divergence · Rotationnel · Laplacien · Propriétés fondamentales}
\end{center}
\noindent\rule{\linewidth}{0.8pt}

\begin{multicols}{2}

\section{Gradient d'un champ scalaire}

\begin{grisbox}
Pour $f(x,y,z)$ scalaire :
$$\boxed{\nabla f = \left(\frac{\partial f}{\partial x},\ \frac{\partial f}{\partial y},\ \frac{\partial f}{\partial z}\right)}$$
C'est un \textbf{vecteur}.
\end{grisbox}

\textbf{Interprétation géométrique :}
\begin{itemize}[leftmargin=1em,itemsep=1pt]
  \item $\nabla f$ est \textbf{perpendiculaire} aux surfaces de niveau $f=\mathrm{cst}$
  \item Il pointe dans la direction de \textbf{plus forte augmentation} de $f$
  \item $\|\nabla f\|$ mesure le taux de variation maximal
\end{itemize}

\section{Divergence d'un champ vectoriel}

\begin{grisbox}
Pour $\mathbf{F}=(P,Q,R)$ champ vectoriel :
$$\boxed{\mathrm{div}\,\mathbf{F} = \frac{\partial P}{\partial x} + \frac{\partial Q}{\partial y} + \frac{\partial R}{\partial z}}$$
C'est un \textbf{scalaire}.
\end{grisbox}

\textbf{Interprétation :} mesure le flux sortant local. Si $\mathrm{div}\,\mathbf{F}>0$, il y a une source ; si $\mathrm{div}\,\mathbf{F}<0$, un puits.

\section{Rotationnel d'un champ vectoriel}

\begin{grisbox}
Pour $\mathbf{F}=(P,Q,R)$ :
$$\boxed{\mathrm{rot}\,\mathbf{F} = \begin{vmatrix}\mathbf{i} & \mathbf{j} & \mathbf{k}\\\partial_x & \partial_y & \partial_z\\ P & Q & R\end{vmatrix}}$$
$$= \!\left(\frac{\partial R}{\partial y}-\frac{\partial Q}{\partial z},\ \frac{\partial P}{\partial z}-\frac{\partial R}{\partial x},\ \frac{\partial Q}{\partial x}-\frac{\partial P}{\partial y}\right)$$
C'est un \textbf{vecteur}.
\end{grisbox}

\textbf{Interprétation :} mesure la rotation locale du champ. Un champ \textbf{irrotationnel} ($\mathrm{rot}\,\mathbf{F}=\mathbf{0}$) ne tourne pas.

\section{Laplacien}

\begin{grisbox}
\textbf{D'un champ scalaire $f$} :
$$\boxed{\Delta f = \frac{\partial^2 f}{\partial x^2} + \frac{\partial^2 f}{\partial y^2} + \frac{\partial^2 f}{\partial z^2}}$$
C'est un \textbf{scalaire}. On peut aussi écrire $\Delta f = \mathrm{div}(\nabla f)$.

\textbf{Fonction harmonique :} $f$ telle que $\Delta f = 0$.
\end{grisbox}

\section{Propriétés fondamentales}

\begin{theobox}
\textbf{Théorème 1.} Pour tout champ gradient $\mathbf{F}=\nabla f$ :
$$\mathrm{rot}(\nabla f) = \mathbf{0}$$
Tout champ gradient est \textbf{irrotationnel}.

\vspace{3pt}
\textbf{Théorème 2.} Pour tout champ vectoriel $\mathbf{F}$ :
$$\mathrm{div}(\mathrm{rot}\,\mathbf{F}) = 0$$
La divergence d'un rotationnel est nulle.
\end{theobox}

\section{Applications géométriques}

\textbf{Vecteur normal à $F(x,y,z)=0$ en $M_0$} : $\nabla F(M_0)$

\textbf{Plan tangent à $F(x,y,z)=0$ en $M_0=(x_0,y_0,z_0)$} :
$$\frac{\partial F}{\partial x}(M_0)(x-x_0) + \frac{\partial F}{\partial y}(M_0)(y-y_0) + \frac{\partial F}{\partial z}(M_0)(z-z_0) = 0$$

\section{Tableau récapitulatif}

\begin{grisbox}
\renewcommand{\arraystretch}{1.3}
\begin{tabular}{@{}lll@{}}
\textbf{Opérateur} & \textbf{Entrée} & \textbf{Sortie}\\
\hline
$\nabla f$ (gradient) & scalaire & vecteur \\
$\mathrm{div}\,\mathbf{F}$ & vecteur & scalaire \\
$\mathrm{rot}\,\mathbf{F}$ & vecteur & vecteur \\
$\Delta f$ (laplacien) & scalaire & scalaire \\
\end{tabular}
\end{grisbox}

\begin{alertbox}
\textbf{Erreurs fréquentes :}
\begin{itemize}[leftmargin=1em,itemsep=0pt]
  \item Confondre gradient (vecteur) et divergence (scalaire)
  \item Appliquer le rotationnel à un scalaire
  \item Oublier que $\mathrm{div}(\nabla f) = \Delta f$ (pas $\nabla^2 f$ vectoriel !)
\end{itemize}
\end{alertbox}

\end{multicols}

\end{document}
